\documentclass[11pt]{article}

\usepackage[margin=1in]{geometry}
\usepackage{booktabs}
\usepackage{enumitem}
\usepackage{hyperref}
\usepackage{url}
\usepackage{xcolor}

\hypersetup{
  colorlinks=true,
  linkcolor=blue,
  citecolor=blue,
  urlcolor=blue
}

\title{Memory as a Control System for Production AI Agents}

\author{
  Vorojar\\
  AgentClaw Project\\
  \url{https://github.com/vorojar/AgentClaw}
}

\date{May 2026}

\begin{document}
\maketitle

\begin{abstract}
Long-term memory is often added to language-agent systems as a retrieval problem:
extract facts, store them, retrieve the top-ranked entries, and inject them into
the prompt. This framing is insufficient for production agents because retrieved
memory is not passive evidence once it enters the model context. It becomes part
of the policy that steers the next action. A memory can be true, textually
relevant, and still harmful for the current task. This paper presents a control
system view of agent memory based on an open-source case study, AgentClaw. The
design separates atomic evidence, scene-level memory, and stable profile memory;
selects an Active Memory set before prompt injection; records memory-use
telemetry; provides edit, deprecate, and merge governance primitives; and turns
memory failures into replayable regression scenarios. We describe the
architecture, failure modes, and a deterministic replay suite that covers stale
preferences, conflicting memories, underspecified follow-up turns, layered
aggregation, telemetry, and safe fallback behavior. The contribution is not a
new benchmark result, but an engineering pattern: production agent memory should
be measured by governed influence, not merely by retrieval recall.
\end{abstract}

\section{Introduction}

Language agents are increasingly expected to perform long-running work across
tools, files, browser sessions, calendars, communication channels, and user
preferences. Memory is a natural requirement for such systems. A useful agent
should remember stable user preferences, prior decisions, long-lived project
facts, and task-specific delivery rules.

The common implementation is retrieval-augmented generation (RAG): store memory
entries, retrieve entries relevant to the current query, and inject them into the
prompt. RAG is useful for grounding language models in external knowledge
\cite{lewis2020rag}. However, agent memory differs from document retrieval in a
critical way. A retrieved memory does not only answer a question. It can change
the agent's behavior.

Consider a real production-style failure pattern. A user previously asked for
dark, technical presentation decks. Later, the user asks for a business
sponsorship deck where a bright, commercial style is more appropriate. If the old
preference is retrieved and placed into the prompt, the model is likely to obey
it. The memory was true, but it should not have influenced this task.

Now consider an underspecified follow-up. The user first asks for a presentation
file. One turn later, the user asks, ``what should be delivered in the end?'' A
memory selector that only sees the latest sentence may miss that the task is
still a presentation task. If the trace contains shared phrases, an unrelated
profile memory can be retrieved by lexical overlap. This is not simply a bad
embedding. It is a control failure: the system allowed the wrong memory to
influence the next action.

This paper argues that production agent memory should be treated as a control
surface. A memory should influence an agent only when the runtime can explain why
it was recalled, record that it was used, test whether it helped, and remove or
supersede it when it causes harm.

We make four contributions:
\begin{enumerate}[leftmargin=*]
  \item We define memory-induced policy failures for production language agents,
  including stale preference injection, conflicting memories, missing turn
  context, stale embeddings after edit, and unsafe selector fallback.
  \item We present a layered memory architecture with Active Memory selection,
  deterministic fallback, telemetry, and governance primitives.
  \item We describe a trace-replay method for converting memory failures into
  deterministic regression tests.
  \item We report a case study from AgentClaw, an open-source multi-channel agent
  workbench, with replay scenarios and source-level artifacts.
\end{enumerate}

\section{Background and Related Work}

RAG combines parametric model knowledge with non-parametric retrieval
\cite{lewis2020rag}. This pattern has become a default way to attach external
knowledge to language models. Tool-using and interactive agents, such as ReAct,
interleave reasoning with actions in external environments \cite{yao2023react}.
Systems such as Reflexion use verbal feedback to improve later attempts
\cite{shinn2023reflexion}, while Generative Agents use memory retrieval,
reflection, and planning to simulate persistent behavior \cite{park2023generative}.
MemGPT frames memory management as an operating-system-like problem over limited
context \cite{packer2023memgpt}. Long-context work also shows that merely placing
information in a large context does not guarantee reliable use
\cite{liu2023lost}.

These lines of work motivate persistent agent state, but production systems need
an additional operational layer. The question is not only which memory is
similar to the current query. The question is whether a memory should be allowed
to influence a concrete decision under current task constraints. That requires
provenance, selection, telemetry, replay, and cleanup.

\section{Failure Model}

We distinguish document relevance from decision relevance. A memory entry is
document-relevant when its text is semantically close to the current request. It
is decision-relevant when it should influence the next agent action. Production
failures occur when document relevance is mistaken for decision relevance.

\begin{table}[h]
\centering
\small
\begin{tabular}{lll}
\toprule
Failure & Example & Root cause \\
\midrule
Stale preference & Old dark-slide preference affects a business deck & No scope or lifecycle model \\
Conflict & Dark style and white-blue style both enter prompt & Superseded memories remain active \\
Missing context & Follow-up omits ``PPTX'' & Selector sees only current turn \\
Stale embedding & Edited memory still retrieves as old content & Edit does not refresh vector state \\
Unsafe fallback & Selector fails, all candidates are injected & Fallback expands rather than reduces risk \\
No replay & Prompt patch fixes one wording only & Failure is not made executable \\
\bottomrule
\end{tabular}
\caption{Memory-induced policy failures in production agents.}
\label{tab:failures}
\end{table}

The central observation is that prompt-injected memory becomes policy-like
material. It changes the model's action distribution. Therefore, memory runtime
design should resemble a feedback control system more than a passive key-value
store.

\section{System Design}

\subsection{Layered Memory}

AgentClaw separates long-term memory into three layers:

\begin{table}[h]
\centering
\small
\begin{tabular}{llll}
\toprule
Layer & Meaning & Example & Runtime role \\
\midrule
L1 & Atomic evidence & User said previews come before final files & Auditable source fact \\
L2 & Scene memory & In PPTX delivery, preview first, then send .pptx & Task-specific rule \\
L3 & Stable profile & User values validation before delivery & Cross-task preference \\
\bottomrule
\end{tabular}
\caption{Layered memory separates evidence from reusable behavioral summaries.}
\label{tab:layers}
\end{table}

L1 memories preserve traceable evidence. L2 and L3 memories compact repeated
evidence into higher-level behavior with provenance links back to source memory
identifiers. A typical L2 memory records a scene name, confidence, source memory
IDs, and status. This lets the agent consume compact behavioral guidance while
maintaining an audit trail.

\subsection{Active Memory Selection}

The Active Memory stage runs before prompt injection. It searches broadly, ranks
candidates deterministically, then chooses a small set of entries that are
directly useful for the current task.

\begin{figure}[h]
\centering
\fbox{\begin{minipage}{0.92\linewidth}
\small
\textbf{Control loop.}
Write memory with provenance $\rightarrow$ aggregate into L1/L2/L3 layers
$\rightarrow$ rank candidates using current and recent user turns
$\rightarrow$ select Active Memory
$\rightarrow$ inject only selected memories
$\rightarrow$ record \texttt{memory\_usage}
$\rightarrow$ replay outcomes
$\rightarrow$ edit, deprecate, merge, or retire harmful memories.
\end{minipage}}
\caption{A governed memory loop treats recall as measurable influence.}
\label{fig:control}
\end{figure}

The query is built from the current user input and recent user turns. This
protects continuity when a follow-up omits the original task noun. The selector
is allowed to return no entries. If the provider-backed selector fails, the
system uses a deterministic fallback over the top-ranked candidates; it never
injects all candidates. The fallback invariant is:

\begin{quote}
When selection confidence is low, reduce prompt influence rather than widen it.
\end{quote}

\subsection{Governance Primitives}

Memory governance is part of the runtime, not only an admin interface.

\begin{itemize}[leftmargin=*]
  \item \textbf{Edit}: when memory content changes, retrieval state such as
  embeddings must be regenerated.
  \item \textbf{Deprecate}: stale or harmful memories are excluded from recall
  without deleting their audit trail.
  \item \textbf{Merge}: fragmented memories are superseded by a canonical memory
  with provenance links.
  \item \textbf{Janitor}: repeated polluting behavior can trigger review or
  retirement.
\end{itemize}

\subsection{Telemetry}

Every prompt-injected memory should be observable. AgentClaw records memory-use
rows with conversation ID, memory ID, source, and selection path. This supports
three questions: which memories influenced a response, whether the relevant
memory was selected, and whether an irrelevant memory polluted the prompt.

\section{Trace Replay}

Component tests are insufficient for memory failures because the bad behavior
emerges from the interaction of conversation state, memory state, selector
behavior, prompt assembly, and tool execution. Trace replay captures the minimum
state needed to recreate the decision boundary.

A replay fixture should include:
\begin{enumerate}[leftmargin=*]
  \item the current user utterance;
  \item recent user turns that define the task;
  \item memory rows that could influence selection;
  \item model or provider behavior relevant to the selector;
  \item the expected user-visible outcome or prompt-injection invariant.
\end{enumerate}

For memory selection, the assertion should not be ``some memory was returned.''
It should verify that the right memory was injected and that a plausible
polluting memory was not injected.

\section{Case Study: AgentClaw}

AgentClaw is an open-source TypeScript monorepo for a multi-channel personal
agent and agent-hosting platform. It includes persistent conversation storage,
long-term memory, tools, skills, traces, telemetry, and multiple user channels.
The memory system is implemented primarily in the \texttt{@agentclaw/memory} and
\texttt{@agentclaw/core} packages.

The case-study scenario comes from presentation delivery. The system seeds an old
preference stating that the user prefers dark PPTX backgrounds. A later turn says
the user now wants white backgrounds with blue emphasis, and requires preview
acceptance before final delivery. The expected behavior is that the old
preference is deprecated, the new preference supersedes it, L2/L3 layered memory
is created, active recall injects useful memory, and telemetry records the use.

\begin{table}[h]
\centering
\small
\begin{tabular}{ll}
\toprule
Scenario & Expected invariant \\
\midrule
Conflict deprecates old memory & Old memory status becomes deprecated \\
Deprecated hidden from search & Search excludes deprecated memory \\
New memory supersedes old & New memory metadata references old memory \\
L2 scene aggregate & Scene memory is created or updated \\
L2 evidence chain & Scene memory has source memory IDs \\
L3 profile aggregate & Stable profile is created or updated \\
L3 evidence chain & Profile memory has source memory IDs \\
Layered prompt injection & Prompt includes L2 or L3 memory \\
Memory-use telemetry & \texttt{memory\_usage} receives rows \\
Safe offload canvas & Large tool output is externalized safely \\
\bottomrule
\end{tabular}
\caption{Deterministic replay checks in \texttt{scripts/memory-10-case-eval.mts}.}
\label{tab:replay}
\end{table}

A separate scenario replay seeds two memories that share a trace phrase: a
relevant PPTX delivery memory and an irrelevant food preference. The current user
turn asks what final file should be delivered but does not repeat ``PPTX.'' The
test passes only if the active prompt contains the PPTX delivery rule and omits
the unrelated preference.

\section{Discussion}

\subsection{Why This Is Not Just Better Retrieval}

Better retrieval can reduce false positives, but it does not solve memory
lifecycle. A stale preference can be highly retrievable. A conflicting memory can
be semantically relevant. A current-turn query can be under-specified. A selector
can fail. Production memory needs a policy for each of these cases.

\subsection{Evaluation Limits}

This paper reports deterministic engineering evidence rather than a broad
benchmark across unrelated agents. The scenarios are intentionally small and
causal. They are designed to fail when a specific runtime invariant regresses.
The next step is to publish a larger benchmark of memory pollution, conflict
resolution, and follow-up continuity across tasks and models.

\subsection{Operational Risks}

Memory governance introduces product decisions. Deprecation can hide useful old
context. Aggressive selection can omit subtle user preferences. Telemetry can
store sensitive references if not designed carefully. Production systems should
combine selection with user-visible memory management, namespace isolation,
sensitive-content scanning, and retention policies.

\section{Submission and Artifact Availability}

The implementation and replay scripts are available in the AgentClaw repository:
\url{https://github.com/vorojar/AgentClaw}. The technical essay that motivated
this paper is available at
\url{https://agentengineering.bibidu.com/blog/memory-control-system}.

\section{Conclusion}

Agent memory is not only a storage or retrieval problem. It is a control problem
over what information is allowed to influence the next decision. Production
agents need layered memory, Active Memory selection, deterministic fallback,
telemetry, governance, and trace replay. The practical standard is simple: a
memory should not enter the prompt unless the system can justify, observe, test,
and retire its influence.

\bibliographystyle{plain}
\bibliography{references}

\end{document}
